\uuid{yvJc}
\exo7id{4245}
\titre{exo7 4245}
\auteur{quercia}
\organisation{exo7}
\datecreate{2010-03-12}
\isIndication{false}
\isCorrection{false}
\chapitre{Calcul d'intégrales}
\sousChapitre{Longueur, aire, volume}
\cpge{mpsi-mp2i.integration.c, pcsi.integration.b, ptsi.integration.b, mpsi-mp2i.integration.d, pcsi.integration.c, ptsi.integration.c, bcpst1.integration-segment.b}
\module{Analyse}
\niveau{L1}
\difficulte{}

\contenu{
\texte{
Soit $f : {[a,b]} \to {\R}$ de classe $\mathcal{C}^1$ telle que $f(a) = f(b) = 0$.
On pose $M' = \|f'\|_\infty$.
}
\begin{enumerate}
    \item \question{En majorant $f$ par une fonction affine par morceaux, démontrer que
      $\left| \int_{t=a}^b f(t)\,d t \right| \le M'\frac {(b-a)^2}4$.}
    \item \question{Quand y a-t-il égalité ?}
\end{enumerate}
}
